\documentclass[11pt]{article}
% The preamble every project's paper shares: % The preamble every project's paper shares: % The preamble every project's paper shares: \input{../../preamble} from
% projects/<project>/paper/. Boards come from chessboard and xskak, plots from
% pgfplots; every number a paper prints comes from its generated.tex.
\usepackage[LSBC4,T1]{fontenc}
\usepackage[utf8]{inputenc}
\usepackage{lmodern}
\usepackage[margin=1in]{geometry}
\usepackage{microtype}
\usepackage{amsmath}
\usepackage{booktabs}
\usepackage{tabularx}
\usepackage{array}
\usepackage{xcolor}
\usepackage[obeyspaces,spaces]{url}
\usepackage{xskak}
\usepackage{pgfplots}
\pgfplotsset{compat=1.18}
\usepackage{caption}
\usepackage{subcaption}
\usepackage{enumitem}
\usepackage[section]{placeins}
\usepackage[hidelinks]{hyperref}

\setchessboard{boardfontencoding=LSBC4, boardfontfamily=skaknew, blackfieldcolor=black!22,
  whitefieldcolor=white, setfontcolors, showmover=false, label=true, labelfontsize=6pt,
  boardfontsize=14pt}

% A move or line in UCI, and a position in FEN, both breakable at a slash.
\newcommand{\uci}[1]{\texttt{#1}}
\newcommand{\fen}[1]{{\footnotesize\path{#1}}}
\newcolumntype{L}{>{\raggedright\arraybackslash}X}

\setlength{\parskip}{0.4em}
\setlength{\parindent}{0pt}
 from
% projects/<project>/paper/. Boards come from chessboard and xskak, plots from
% pgfplots; every number a paper prints comes from its generated.tex.
\usepackage[LSBC4,T1]{fontenc}
\usepackage[utf8]{inputenc}
\usepackage{lmodern}
\usepackage[margin=1in]{geometry}
\usepackage{microtype}
\usepackage{amsmath}
\usepackage{booktabs}
\usepackage{tabularx}
\usepackage{array}
\usepackage{xcolor}
\usepackage[obeyspaces,spaces]{url}
\usepackage{xskak}
\usepackage{pgfplots}
\pgfplotsset{compat=1.18}
\usepackage{caption}
\usepackage{subcaption}
\usepackage{enumitem}
\usepackage[section]{placeins}
\usepackage[hidelinks]{hyperref}

\setchessboard{boardfontencoding=LSBC4, boardfontfamily=skaknew, blackfieldcolor=black!22,
  whitefieldcolor=white, setfontcolors, showmover=false, label=true, labelfontsize=6pt,
  boardfontsize=14pt}

% A move or line in UCI, and a position in FEN, both breakable at a slash.
\newcommand{\uci}[1]{\texttt{#1}}
\newcommand{\fen}[1]{{\footnotesize\path{#1}}}
\newcolumntype{L}{>{\raggedright\arraybackslash}X}

\setlength{\parskip}{0.4em}
\setlength{\parindent}{0pt}
 from
% projects/<project>/paper/. Boards come from chessboard and xskak, plots from
% pgfplots; every number a paper prints comes from its generated.tex.
\usepackage[LSBC4,T1]{fontenc}
\usepackage[utf8]{inputenc}
\usepackage{lmodern}
\usepackage[margin=1in]{geometry}
\usepackage{microtype}
\usepackage{amsmath}
\usepackage{booktabs}
\usepackage{tabularx}
\usepackage{array}
\usepackage{xcolor}
\usepackage[obeyspaces,spaces]{url}
\usepackage{xskak}
\usepackage{pgfplots}
\pgfplotsset{compat=1.18}
\usepackage{caption}
\usepackage{subcaption}
\usepackage{enumitem}
\usepackage[section]{placeins}
\usepackage[hidelinks]{hyperref}

\setchessboard{boardfontencoding=LSBC4, boardfontfamily=skaknew, blackfieldcolor=black!22,
  whitefieldcolor=white, setfontcolors, showmover=false, label=true, labelfontsize=6pt,
  boardfontsize=14pt}

% A move or line in UCI, and a position in FEN, both breakable at a slash.
\newcommand{\uci}[1]{\texttt{#1}}
\newcommand{\fen}[1]{{\footnotesize\path{#1}}}
\newcolumntype{L}{>{\raggedright\arraybackslash}X}

\setlength{\parskip}{0.4em}
\setlength{\parindent}{0pt}

% Written by projects/rules-after-check/paper/generate.mjs; do not edit by hand.

\newcommand{\CheckCommit}{3437a71184}
\newcommand{\ScoreCommit}{f4bd63f19b}
\newcommand{\RecordsCommit}{aed714ada5}
\newcommand{\TrainMonths}{July and August 2026}
\newcommand{\HeldMonth}{September 2026}
\newcommand{\BandLow}{1250}
\newcommand{\BandHigh}{1649}
\newcommand{\HisHigh}{1449}
\newcommand{\Inherited}{127,914}
\newcommand{\NRules}{13}
\newcommand{\Draws}{2,000}
\newcommand{\Seed}{20261018}
\newcommand{\MinHeld}{300}
\newcommand{\RunHours}{4.7}
\newcommand{\WorkerHours}{9.3}
\newcommand{\ControlTrain}{$-$1.0}
\newcommand{\ControlTrainCi}{$-$6.5 to 4.5}
\newcommand{\ControlTrainN}{292}
\newcommand{\ControlTrainGames}{288}
\newcommand{\ControlHeld}{2.3}
\newcommand{\ControlHeldCi}{$-$6.0 to 10.1}
\newcommand{\ControlHeldN}{175}
\newcommand{\ControlHeldGames}{173}
\newcommand{\ControlShareTrain}{0.9}
\newcommand{\ControlShareHeld}{1.1}
\newcommand{\PrimaryTrain}{30,758}
\newcommand{\PrimaryHeld}{15,489}
\newcommand{\ControlTable}{%
\begin{table}[ht]
\centering\footnotesize
\begin{tabular}{@{}lrrrl@{}}
\toprule
Period & \begin{tabular}[b]{@{}l@{}}Quiet positions with \\ a safe rule move\end{tabular} & \begin{tabular}[b]{@{}l@{}}Control \\ positions\end{tabular} & Games & \begin{tabular}[b]{@{}l@{}}$R_s$ minus $R$, \\ points\end{tabular} \\
\midrule
Train & 30,758 & 292 & 288 & $-$1.0 {\scriptsize($-$6.5, 4.5)} \\
Held out & 15,489 & 175 & 173 & 2.3 {\scriptsize($-$6.0, 10.1)} \\
\bottomrule
\end{tabular}
\caption{The control, pooled over the 13 rules: the quiet positions where the rule's move $R$ fails the app's blunder check and some rule move passes it, and the soundness of $R_s$, the rule's move after the check, minus $R$'s, paired by position, with its 95\% interval from 2,000 draws over games. The rule asked for the interval to lie above 0 on both periods.}
\label{tab:control}
\end{table}
}
\newcommand{\NWith}{twelve}
\newcommand{\ThinName}{the outside passed pawn}
\newcommand{\ThinHeld}{12}
\newcommand{\ThinTrain}{17}
\newcommand{\LiftLow}{1.3}
\newcommand{\LiftLowName}{the rule of the square}
\newcommand{\LiftHigh}{32.0}
\newcommand{\LiftHighName}{the pawn break}
\newcommand{\RsRLow}{95.5}
\newcommand{\RsRHigh}{100.0}
\newcommand{\NsNLow}{93.9}
\newcommand{\NsNHigh}{100.0}
\newcommand{\NaturalMax}{0.9}
\newcommand{\ThinSafe}{8}
\newcommand{\LowSafeNames}{cutting the king off, the passed pawn and the worst piece}
\newcommand{\LowSafeLow}{76}
\newcommand{\LowSafeHigh}{82}
\newcommand{\QuietLow}{12}
\newcommand{\QuietHigh}{72}
\newcommand{\QuietLowName}{the outside passed pawn}
\newcommand{\QuietHighName}{the opposition}
\newcommand{\CondLow}{6.7}
\newcommand{\CondLowName}{the rule of the square}
\newcommand{\CondHigh}{36.4}
\newcommand{\CondHighName}{the pawn break}
\newcommand{\NContent}{ten}
\newcommand{\ContentNotNames}{the rook behind the passed pawn and cutting the king off}
\newcommand{\ContentLow}{3.4}
\newcommand{\ContentLowName}{the opposition}
\newcommand{\ContentHigh}{36.8}
\newcommand{\ContentHighName}{trading when ahead}
\newcommand{\RbpContentTrain}{2.1}
\newcommand{\RbpContentTrainCi}{0.2 to 3.8}
\newcommand{\RbpContentHeld}{2.1}
\newcommand{\RbpContentHeldCi}{$-$0.5 to 4.6}
\newcommand{\CutContentHeld}{$-$1.1}
\newcommand{\CutContentHeldCi}{$-$4.9 to 2.8}
\newcommand{\RatioLow}{1.3}
\newcommand{\RatioHigh}{2.4}
\newcommand{\RatioLowName}{the opposition}
\newcommand{\RatioHighName}{the active king}
\newcommand{\AllLow}{2.1}
\newcommand{\AllHigh}{47.6}
\newcommand{\FreshSearches}{84,198}
\newcommand{\DescriptionTable}{%
\begin{table}[ht]
\centering\scriptsize
\setlength{\tabcolsep}{3pt}
\begin{tabular}{@{}lrrrrrlrl@{}}
\toprule
 & \begin{tabular}[b]{@{}l@{}}Held-out \\ positions\end{tabular} & \begin{tabular}[b]{@{}l@{}}Quiet \\ share\end{tabular} & \begin{tabular}[b]{@{}l@{}}$R_s$ \\ sound\end{tabular} & \begin{tabular}[b]{@{}l@{}}$N_s$ \\ sound\end{tabular} & \begin{tabular}[b]{@{}l@{}}Lift, \\ train\end{tabular} & \begin{tabular}[b]{@{}l@{}}Lift, \\ held out\end{tabular} & \begin{tabular}[b]{@{}l@{}}Lift, all \\ positions\end{tabular} & \begin{tabular}[b]{@{}l@{}}Content, \\ held out\end{tabular} \\
\midrule
Active king & 2,245 & 59 & 72.5 & 82.3 & $-$10.7 & $-$9.8 {\scriptsize($-$11.5, $-$8.2)} & $-$23.3 & 7.5 {\scriptsize(5.7, 9.4)} \\
Opposition & 1,015 & 72 & 85.8 & 94.3 & $-$6.0 & $-$8.5 {\scriptsize($-$10.6, $-$6.4)} & $-$10.6 & 3.4 {\scriptsize(1.3, 5.7)} \\
Rule of the square & 1,552 & 60 & 97.2 & 98.5 & $-$1.2 & $-$1.3 {\scriptsize($-$2.0, $-$0.6)} & $-$2.1 & 10.0 {\scriptsize(8.3, 11.6)} \\
Rook behind the passed pawn & 1,307 & 35 & 51.0 & 75.7 & $-$25.3 & $-$24.6 {\scriptsize($-$27.5, $-$21.6)} & $-$41.3 & 2.1 {\scriptsize($-$0.5, 4.6)} \\
Cutting the king off & 552 & 52 & 53.3 & 82.2 & $-$28.6 & $-$29.0 {\scriptsize($-$33.1, $-$25.0)} & $-$37.6 & $-$1.1 {\scriptsize($-$4.9, 2.8)} \\
Outside passed pawn & 12 & 12 & 58.3 & 75.0 & 0.0 & $-$16.7 {\scriptsize($-$42.9, 0.0)} & $-$4.6 & 8.3 {\scriptsize($-$36.4, 50.0)} \\
Trading when ahead & 1,033 & 27 & 77.6 & 84.4 & $-$6.9 & $-$6.8 {\scriptsize($-$8.8, $-$4.8)} & $-$11.8 & 36.8 {\scriptsize(33.3, 40.2)} \\
Open file & 1,454 & 38 & 54.1 & 71.1 & $-$18.0 & $-$17.0 {\scriptsize($-$20.0, $-$14.0)} & $-$30.5 & 7.0 {\scriptsize(4.3, 9.6)} \\
Rook on the seventh & 1,176 & 31 & 59.5 & 76.7 & $-$16.2 & $-$17.2 {\scriptsize($-$20.1, $-$14.2)} & $-$28.1 & 14.9 {\scriptsize(12.0, 18.0)} \\
Outpost & 1,259 & 34 & 53.8 & 73.9 & $-$20.7 & $-$20.2 {\scriptsize($-$23.0, $-$17.3)} & $-$30.3 & 13.0 {\scriptsize(10.0, 16.0)} \\
Pawn break & 1,620 & 44 & 39.5 & 71.5 & $-$30.5 & $-$32.0 {\scriptsize($-$34.8, $-$29.1)} & $-$47.6 & 3.6 {\scriptsize(0.9, 6.4)} \\
Passed pawn & 1,031 & 32 & 64.3 & 79.1 & $-$14.4 & $-$14.8 {\scriptsize($-$17.5, $-$12.2)} & $-$29.2 & 12.6 {\scriptsize(9.6, 15.6)} \\
Worst piece & 1,233 & 37 & 52.6 & 75.1 & $-$20.1 & $-$22.5 {\scriptsize($-$25.5, $-$19.4)} & $-$36.4 & 4.7 {\scriptsize(2.0, 7.4)} \\
\bottomrule
\end{tabular}
\caption{What the measures describe, with no verdict. Held-out positions: quiet ones where some rule move passes the check. Quiet share: of all the rule's held-out positions, in percent. $R_s$ and $N_s$ sound: the share of those positions where the rule's move and the natural move after the check are sound, in percent. The lift is $R_s$'s soundness minus $N_s$'s in points, with its 95\% interval from 2,000 draws over games (for every rule with data the train interval lies below 0 as well); over all positions it is the same lift over every position with a safe rule move, quiet or not. The content is $R_s$'s soundness minus that of $R_{0s}$, the move Maia-2 likes best among as many safe moves drawn at random.}
\label{tab:descriptions}
\end{table}
}
\newcommand{\NForcing}{eight}
\newcommand{\ForcingAcrossNames}{the open file and the rook on the seventh}
\newcommand{\ForcingMin}{50}
\newcommand{\NBigForcing}{four}
\newcommand{\OfForcing}{$-$4.1}
\newcommand{\OfQuiet}{$-$17.9}
\newcommand{\NPhase}{eight}
\newcommand{\EndgameOnlyNames}{the active king, the opposition, the rule of the square and cutting the king off}
\newcommand{\PpMiddle}{$-$25.1}
\newcommand{\PpEnd}{$-$9.7}
\newcommand{\NMaterial}{eleven}
\newcommand{\NLevel}{nine}
\newcommand{\LevelNotNames}{the rule of the square and cutting the king off}
\newcommand{\SplitTable}{%
\begin{table}[ht]
\centering\scriptsize
\setlength{\tabcolsep}{3pt}
\begin{tabular}{@{}lrrrrrrr@{}}
\toprule
 & \begin{tabular}[b]{@{}l@{}}$R_s$ \\ quiet\end{tabular} & \begin{tabular}[b]{@{}l@{}}$R_s$ \\ forcing\end{tabular} & \begin{tabular}[b]{@{}l@{}}Middle- \\ game\end{tabular} & Endgame & Ahead & Level & Behind \\
\midrule
Active king & $-$10 {\scriptsize(2,244)} & 0 {\scriptsize(1)} & -- & $-$10 {\scriptsize(2,245)} & $-$9 {\scriptsize(778)} & $-$14 {\scriptsize(875)} & $-$4 {\scriptsize(592)} \\
Opposition & $-$8 {\scriptsize(1,015)} & -- & -- & $-$8 {\scriptsize(1,015)} & $-$4 {\scriptsize(230)} & $-$11 {\scriptsize(642)} & $-$6 {\scriptsize(143)} \\
Rule of the square & $-$1 {\scriptsize(1,549)} & 0 {\scriptsize(3)} & -- & $-$1 {\scriptsize(1,552)} & $-$1 {\scriptsize(1,205)} & $-$1 {\scriptsize(329)} & $-$6 {\scriptsize(18)} \\
Rook behind the passed pawn & $-$25 {\scriptsize(1,298)} & 11 {\scriptsize(9)} & $-$33 {\scriptsize(397)} & $-$21 {\scriptsize(910)} & $-$23 {\scriptsize(587)} & $-$31 {\scriptsize(454)} & $-$18 {\scriptsize(266)} \\
Cutting the king off & $-$29 {\scriptsize(552)} & -- & -- & $-$29 {\scriptsize(552)} & $-$16 {\scriptsize(176)} & $-$34 {\scriptsize(341)} & $-$46 {\scriptsize(35)} \\
Trading when ahead & -- & $-$7 {\scriptsize(1,033)} & $-$8 {\scriptsize(806)} & $-$4 {\scriptsize(227)} & $-$7 {\scriptsize(1,033)} & -- & -- \\
Open file & $-$18 {\scriptsize(1,356)} & $-$4 {\scriptsize(98)} & $-$19 {\scriptsize(976)} & $-$12 {\scriptsize(478)} & $-$15 {\scriptsize(351)} & $-$19 {\scriptsize(794)} & $-$14 {\scriptsize(309)} \\
Rook on the seventh & $-$18 {\scriptsize(1,113)} & 2 {\scriptsize(63)} & $-$19 {\scriptsize(347)} & $-$16 {\scriptsize(829)} & $-$12 {\scriptsize(420)} & $-$22 {\scriptsize(476)} & $-$15 {\scriptsize(280)} \\
Outpost & $-$21 {\scriptsize(1,142)} & $-$10 {\scriptsize(117)} & $-$23 {\scriptsize(720)} & $-$17 {\scriptsize(539)} & $-$19 {\scriptsize(413)} & $-$22 {\scriptsize(662)} & $-$16 {\scriptsize(184)} \\
Pawn break & $-$32 {\scriptsize(1,619)} & 0 {\scriptsize(1)} & $-$34 {\scriptsize(1,239)} & $-$25 {\scriptsize(381)} & $-$29 {\scriptsize(280)} & $-$34 {\scriptsize(1,086)} & $-$24 {\scriptsize(254)} \\
Passed pawn & $-$16 {\scriptsize(870)} & $-$9 {\scriptsize(161)} & $-$25 {\scriptsize(342)} & $-$10 {\scriptsize(689)} & $-$13 {\scriptsize(476)} & $-$17 {\scriptsize(402)} & $-$14 {\scriptsize(153)} \\
Worst piece & $-$23 {\scriptsize(1,210)} & $-$9 {\scriptsize(23)} & $-$25 {\scriptsize(714)} & $-$18 {\scriptsize(519)} & $-$18 {\scriptsize(312)} & $-$27 {\scriptsize(662)} & $-$15 {\scriptsize(259)} \\
\bottomrule
\end{tabular}
\caption{The lift in quiet positions in each cell of a split, held out, in points, with the cell's positions in brackets; a dash is a cell with none. Forcing: $R_s$ is a capture or gives check. The phase is the app's divider's, and material is the mover's, both read before the move by the app at the check's commit. The splits describe and decide nothing.}
\label{tab:splits}
\end{table}
}


\title{Does a rule of thumb choose better once the tactics are checked?\\[0.3em]
\large The app's rules of thumb after its blunder check, in quiet positions from \Inherited{}
Lichess positions refereed by Stockfish}
\author{Chess Coach research lab}
\date{}

\begin{document}
\maketitle

\begin{abstract}
A coach puts the checks, captures and threats first and a rule of thumb second, and so do several of
the app's lessons. The previous study found that no rule of thumb chooses a better move than the one
a player at his level would naturally play, mostly because the rule's move left a piece to be taken
or passed over a tactic. This one tests the rule as the lessons teach it: in positions with no tactic
on the board by the app's reading, and among the moves the app's blunder check passes, is the move a
rule points to, as his level would play it, sound as often as the move his level naturally plays
among the same moves? The rule, fixed before any position was read, first asks a control to show
that the check can rescue a rule's move. It does not: where the rule's own move fails the check, the
next rule move that passes it is no sounder, its soundness minus the rule move's \ControlTrain{}
points on the train months and \ControlHeld{} held out, both intervals across 0, so no rule receives
a verdict. The descriptions
are clear all the same. In a quiet position the check rarely touches the rule's move, and that move
is sound less often than the natural one for all \NWith{} rules with data, on both periods: by
\LiftLow{} points for \LiftLowName{} up to \LiftHigh{} for \LiftHighName{} held out. Over all
positions, tactics included, the gap is \RatioLow{} to \RatioHigh{} times its quiet size. The rules
still point to sounder moves than chance, \NContent{} of them on both periods. A rule of thumb tells
a player where to look and explains a move; at his level it does not choose one, even after the
tactics.
\end{abstract}

\section{The question}

The previous study asked each of the app's rules of thumb to choose a move with no tactics checked
first: among the legal moves that show the rule, the one Maia-2 rates most likely at the mover's
band. No rule chose better than the natural move, the one Maia-2 rates most likely of all, and its
held-out examples said why. Most of the rule's unsound moves lose material by the app's own blunder
check, and most of the rest pass over a check, a capture or a promotion that the natural move took.

A coach would not give the rule that way. Several of the app's lessons begin ``No tactics?'', and
after that study its open-file and outpost lessons put the checks, captures and threats first too.
So the fair test of such a lesson is the rule after the check. In positions with no tactic on the
board, and among the moves the app's blunder check passes, is the move a rule points to, as players
at his level would play it, sound as often as the move they would naturally play among those same
moves? If it is sounder, a lesson can teach the rule as the move to look for once nothing is
hanging. If it is no sounder, the rule explains a move and says where to look, but does not choose
one. If it is less sound, a lesson that teaches it should tell him to compare its move with the
others before playing it.

\section{How it was measured}

\paragraph{The positions.} The \Inherited{} positions the previous study scored for its \NRules{}
rules with a verdict, read and never changed: middlegames and endgames from rated 10-minute Lichess
games with no increment, the mover rated \BandLow{} to \BandHigh{} on Lichess (his band and the one
above), with \TrainMonths{} to train on and \HeldMonth{} held out, at most one position a game for
each rule. Each keeps the rule's legal moves, the moves that study compared and the referee's
evidence for them.

\paragraph{The check.} Two readings by the app's knowledge service at commit \texttt{\CheckCommit}.
A position is \emph{quiet}, the lessons' ``no tactics'', when the mover is not in check and neither
side has a mate in one or a winning capture, by the app's description of the position. A move
\emph{passes the check} when the app's safety check calls it safe: the reply cannot mate at once, no
piece of the mover's is newly left to be taken for a gain (the capturing piece of an even trade is
not counted), and the move deserts no defence.

\paragraph{The moves compared.} In each quiet position, $R_s$ is the rule's move after the check:
among the rule's moves that pass it, the one Maia-2 rates most likely at the mover's band. $N_s$,
the natural move after the check, is Maia-2's most likely move among all the moves that pass.
$R_{0s}$ is the move Maia-2 likes best among as many passing moves as $R_s$ had to choose from,
drawn at random: a rule with no information. $R$ and $N$ are the previous study's rule move and
natural move, before the check. Maia-2 is the app's at commit \texttt{\ScoreCommit}, the one that
study measured, and every $R$ and $N$ was reproduced before a position was scored.

\paragraph{The referee.} A move the previous study scored keeps its soundness, and with seven men or
fewer its stored tablebase probe decides. Every other move is scored by Stockfish through the app's
engine service at \texttt{\ScoreCommit}, at its grading budget, against one fresh search for the best
in that position: a move is \emph{sound} when it stays within the app's window of the best. All
\FreshSearches{} fresh searches found the best that study had stored. The scoring took
\WorkerHours{} worker hours on two workers.

\paragraph{The rule, fixed before any position was read.} The \emph{lift} is $R_s$'s soundness
minus $N_s$'s, paired by position over the quiet positions where some rule move passes the check,
with 95\% intervals from \Draws{} bootstrap draws over games (seed \Seed). A control comes first:
pooled over the rules, in the quiet positions where $R$ fails the check and some rule move passes
it, $R_s$'s soundness minus $R$'s must have its interval above 0 on both periods, or no rule gets a
verdict, since a lift this design could not see would read as none. Then a rule holds after the check
when its lift's interval lies above 0 on the train months and on the held-out month, misleads when it
lies below 0 on both, and does no better than the natural move otherwise; with fewer than
\MinHeld{} held-out positions there is not enough data. The \emph{content} after the check,
$R_s$'s soundness minus $R_{0s}$'s, the same measures over every position with a safe rule move,
quiet or not, and the lift in splits of the positions describe and change no verdict.

\section{The control fails}

In the quiet positions with a safe rule move, the rule's own move fails the check only rarely:
\ControlShareTrain\% of them on the train months and \ControlShareHeld\% held out, which leaves the
control \ControlTrainN{} and \ControlHeldN{} positions (Table~\ref{tab:control}). Where it fails, the
next rule move that passes the check is no sounder: $R_s$'s soundness minus $R$'s is
\ControlTrain{} points (\ControlTrainCi) on the train months and \ControlHeld{} (\ControlHeldCi)
held out. By the rule, then, no rule receives a verdict, and everything below describes. This study
does not separate the two ways the control can fail: a quiet move the check flags may be sound after
all, or the next rule move may be no better.

\ControlTable

\section{What the measures describe}

\paragraph{The check changes little in a quiet position.} Where some rule move passes it, the rule's
own move passes it in \RsRLow\% to \RsRHigh\% of the positions, for every rule on both periods, so
$R_s$ is nearly always $R$; the natural move passes it in \NsNLow\% to \NsNHigh\%, and the check moves
the natural move's soundness by at most \NaturalMax{} points for every rule with data. What the check
removes is whole positions: \ThinName{} has a rule move that passes it in only \ThinSafe\% of its
quiet positions held out, which leaves \ThinHeld{}, too few for a verdict even had the control
passed, and \LowSafeNames{} have one in \LowSafeLow\% to \LowSafeHigh\% of theirs. How many of a
rule's positions are quiet at all varies from \QuietLow\% for \QuietLowName{} to \QuietHigh\% for
\QuietHighName.

\paragraph{The rule's move still trails the natural move.} After the check the rule's move is sound
less often than the natural move, for all \NWith{} rules with data, on the train months and on the
held-out month, every interval below 0 (Table~\ref{tab:descriptions}): by \LiftLow{} points for
\LiftLowName{} up to \LiftHigh{} for \LiftHighName{} held out. Where the two moves differ, the rule's
is sound less often by \CondLow{} points for \CondLowName{} up to \CondHigh{} for \CondHighName. This
is the previous study's finding again, in the positions where its two main reasons, a piece left to
be taken and a tactic passed over, are gone by the app's reading.

\paragraph{The rules still carry information.} Against the move Maia-2 likes best among as many safe
moves drawn at random, \NContent{} rules point to sounder moves on both periods, from \ContentLow{}
points for \ContentLowName{} to \ContentHigh{} for \ContentHighName{} held out. The rook behind the
passed pawn does so on the train months, by \RbpContentTrain{} points (\RbpContentTrainCi), but not
held out (\RbpContentHeld, \RbpContentHeldCi), and cutting the king off reads no differently from
chance, \CutContentHeld{} points (\CutContentHeldCi) held out, as in the previous study.

\DescriptionTable

\section{With a tactic on the board}

Over every position with a safe rule move, quiet or not, the gap is larger for each of the \NWith{}
rules on both periods: \RatioLow{} to \RatioHigh{} times its size in the quiet positions, from
\RatioLowName{} to \RatioHighName, and \AllLow{} to \AllHigh{} points held out
(Table~\ref{tab:descriptions}). The quiet positions are among them, so in the rest, where the mover is
in check or either side has a mate in one or a winning capture, the gap is larger still. A tactic on
the board is where a rule's move falls furthest behind the natural one, which is the order the
lessons teach: the tactics first.

\section{Where the gap is smallest}

The splits describe and decide nothing (Table~\ref{tab:splits}). For all \NForcing{} rules whose
move after the check is a capture or a check in some positions and neither in others, on both
periods, the forcing move trails the natural move by less. Of the \NBigForcing{} rules with at least
\ForcingMin{} forcing moves in each period, \ForcingAcrossNames{} do not clearly trail when forcing,
their intervals across 0 on both periods (the open file: \OfForcing{} points forcing and \OfQuiet{}
otherwise, held out). For all \NPhase{} rules that come up in both phases, the gap is larger in the
middlegame than in the endgame on both periods, as with the passed pawn (\PpMiddle{} points against
\PpEnd{} held out); \EndgameOnlyNames{} come up almost only in endgames. And for \NLevel{} of the
\NMaterial{} rules with positions ahead, level and behind, the gap is largest with level material
on both periods; the exceptions are \LevelNotNames. As in the previous study, level material may be
in part the referee's doing, an inference this study does not test: a fixed window in win\% is
narrower in centipawns where the game is balanced.

\SplitTable

\section{What it means for the app}

Nothing is ported, since there is no verdict. The descriptions agree with the previous study's
reading and narrow it. Once the checks, captures and threats are seen, the move a rule of thumb
points to is still sound less often than the move his level plays naturally, and most of all with a
tactic on the board. A rule tells him where to look and explains a move; it does not choose one.

That is how the app now teaches its rules. Its worst-piece lesson begins ``No tactics?'', its
pawn-break lesson says to check what a break leaves hanging before playing it, and since chess PR
165, after the commits measured here, its open-file and outpost lessons start from the checks,
captures and threats, the open-file lesson saying that the file tells him where to look, not what to
play. Game Review already uses the rules that way, as the previous study found.

A round that wants a verdict needs a control with more to see. The check changes the rule's move in
\ControlShareTrain\% and \ControlShareHeld\% of the quiet positions with a safe rule move, so its
rescues are few; a control built from positions where a rule's move is known to leave a piece to be
taken while another rule move is safe would test the check itself.

\section{Caveats}

\begin{itemize}[nosep]
\item \textbf{No verdict.} The control failed, so every number here describes; none was a test the
  rule set.
\item \textbf{Not his games.} The positions are Lichess players' in his band and the band above;
  his own games are too few for most rules.
\item \textbf{Quiet by board code.} No mate in one and no winning capture for either side is the
  app's reading of ``no tactics''. A deeper tactic, a fork two moves away, still counts as quiet.
\item \textbf{One move for a plan.} A rule is tested as the next move his level would play for it,
  not as a plan over several moves, and a sound move can still be the second-best plan.
\item \textbf{A strong baseline.} The natural move is Maia-2's most likely move at the mover's band,
  which the previous study found sound more often than the move actually played.
\item \textbf{Two app commits.} The check is the app's reading at \texttt{\CheckCommit}, after a fix
  to its even-trade reading; Maia-2 and the referee are the previous study's, at
  \texttt{\ScoreCommit}, so every $R$ and $N$ is the one it scored.
\item \textbf{No examples.} The rule draws held-out examples only for rules with a verdict, so this
  paper shows no positions.
\end{itemize}

\end{document}
